\chapter{مسئله تناظر پست و کاهش نگاشتی}

تصمیم‌ناپذیری فقط به مسئله‌هایی درباره ماشین‌ها محدود نیست. در این فصل یک مسئله ساده درباره رشته‌ها می‌بینیم که تصمیم‌ناپذیر است: \term{مسئله تناظر پست}{Post Correspondence Problem} یا \lr{PCP}. سپس مفهوم کاهش را با تعریف دقیق \textbf{کاهش نگاشتی} صوری می‌کنیم.

\section{مسئله تناظر پست}

یک \textbf{دومینو} دو نیمه دارد که روی هر کدام یک رشته نوشته شده، مثل $\dom{a}{ab}$. یک \textbf{مجموعه دومینو} داریم، مثلاً
\[\left\{\dom{b}{ca},\ \dom{a}{ab},\ \dom{ca}{a},\ \dom{abc}{c}\right\}.\]
هدف پیدا کردن یک \textbf{تطابق} است: فهرستی از دومینوهای این مجموعه (با تکرار مجاز) که اگر کنار هم بچینیم، رشته بالا با رشته پایین یکی شود. برای مجموعه بالا:
\[\dom{a}{ab}\dom{b}{ca}\dom{ca}{a}\dom{a}{ab}\dom{abc}{c}:\qquad \text{\rl{بالا}}=abcaaabc=\text{\rl{پایین}}.\]
بعضی مجموعه‌ها تطابق ندارند. مثلاً در $\left\{\dom{abc}{ab},\ \dom{ca}{a},\ \dom{acc}{ba}\right\}$ رشته بالای هر دومینو از رشته پایینش بلندتر است، پس رشته بالا همیشه بلندتر از پایین می‌ماند.

\begin{definition}[\lr{PCP}]
یک نمونه از \lr{PCP} مجموعه‌ای از دومینوهاست:
\[P=\left\{\left[\frac{t_1}{b_1}\right],\left[\frac{t_2}{b_2}\right],\dots,\left[\frac{t_k}{b_k}\right]\right\}.\]
یک \textbf{تطابق} دنباله‌ای $i_1,i_2,\dots,i_l$ (با $l\ge1$) است که $t_{i_1}t_{i_2}\cdots t_{i_l}=b_{i_1}b_{i_2}\cdots b_{i_l}$. مسئله این است که تعیین کنیم $P$ تطابق دارد یا نه:
\[\prob{PCP}=\{\enc{P}\st P\ \text{\rl{یک نمونه \lr{PCP} با یک تطابق است}}\}.\]
\end{definition}

\begin{theorem}
$\prob{PCP}$ تصمیم‌ناپذیر است.
\end{theorem}

\begin{proofidea}
از $\ATM$ با تاریخچه‌های محاسبه پذیرنده کاهش می‌دهیم. از هر ماشین $M$ و ورودی $w$ یک نمونه $P$ می‌سازیم که تطابق‌هایش دقیقاً تاریخچه‌های محاسبه پذیرنده $M$ روی $w$ هستند. یعنی ساختن یک تطابق، همان \textbf{شبیه‌سازی} $M$ روی $w$ است. پس اگر بتوانیم تشخیص دهیم تطابقی هست، می‌فهمیم $M$ رشته $w$ را می‌پذیرد یا نه.

سه جزئیات فنی:
\begin{itemize}
\item فرض می‌کنیم $M$ روی $w$ هرگز سعی نمی‌کند هد را از انتهای چپ نوار بیرون ببرد ($M$ را می‌توان طوری تغییر داد که این‌طور باشد).
\item اگر $w=\eps$، به جای آن از $\blank$ استفاده می‌کنیم.
\item ابتدا مسئله را برای حالتی حل می‌کنیم که تطابق باید \textbf{با دومینوی اول} $\left[\frac{t_1}{b_1}\right]$ شروع شود. به این مسئله \textbf{\lr{PCP} اصلاح‌شده} (\lr{MPCP}) می‌گوییم:
\[\prob{MPCP}=\{\enc{P}\st P\ \text{\rl{نمونه‌ای از \lr{PCP} با تطابقی است که با دومینوی اول شروع می‌شود}}\}.\]
\end{itemize}
\end{proofidea}

\begin{proof}
فرض کنید ماشین $R$ زبان $\prob{PCP}$ را تصمیم بگیرد. ماشین $S$ را می‌سازیم که $\ATM$ را تصمیم بگیرد. فرض کنید $M=(Q,\Sigma,\Gamma,\delta,q_0,\qacc,\qrej)$. ماشین $S$ یک نمونه $P$ از \lr{PCP} می‌سازد که تطابق دارد اگر و تنها اگر $M$ رشته $w$ را بپذیرد. برای این کار، $S$ ابتدا یک نمونه $P'$ از \lr{MPCP} در هفت بخش می‌سازد.

\textbf{بخش ۱ (پیکربندی اول):} دومینوی زیر را به عنوان دومینوی اول در $P'$ بگذار:
\[\dom{\#}{\#q_0w_1w_2\cdots w_n\#}.\]
چون تطابق باید با این دومینو شروع شود، رشته پایین با $C_1$ (پیکربندی شروع) شروع می‌شود. حالا رشته پایین از بالا جلوتر است. برای اینکه بالا به پایین برسد باید دومینوهایی اضافه کنیم که بالایشان $C_1$ باشد؛ این دومینوها طوری ساخته شده‌اند که پایینشان $C_2$ (پیکربندی بعدی) را تولید کند. پس تطابق مجبور است گام به گام محاسبه $M$ را شبیه‌سازی کند.

\textbf{بخش ۲ (حرکت هد به راست):} برای هر $a,b\in\Gamma$ و هر $q,r\in Q$ با $q\neq\qrej$:
\[\text{\rl{اگر}}\ \delta(q,a)=(r,b,R)\ \text{\rl{دومینوی}}\ \dom{qa}{br}\ \text{\rl{را در $P'$ بگذار.}}\]

\textbf{بخش ۳ (حرکت هد به چپ):} برای هر $a,b,c\in\Gamma$ و هر $q,r\in Q$ با $q\neq\qrej$:
\[\text{\rl{اگر}}\ \delta(q,a)=(r,b,L)\ \text{\rl{دومینوی}}\ \dom{cqa}{rcb}\ \text{\rl{را در $P'$ بگذار.}}\]

\textbf{بخش ۴ (نمادهای دور از هد):} برای هر $a\in\Gamma$ دومینوی $\dom{a}{a}$ را بگذار. این دومینوها بخش‌هایی از نوار را که تغییر نمی‌کنند کپی می‌کنند.

\textbf{بخش ۵ (جداکننده و نوار نامحدود):} دومینوهای $\dom{\#}{\#}$ و $\dom{\#}{\sqcup\#}$ را بگذار. اولی جداکننده پیکربندی‌ها را کپی می‌کند و دومی یک خانه خالی به انتهای پیکربندی اضافه می‌کند تا نوار نامحدود شبیه‌سازی شود (وقتی هد به انتهای راست می‌رسد).

\begin{example}
فرض کنید $\Gamma=\{0,1,2,\blank\}$، $w=0100$، حالت شروع $q_0$ و $\delta(q_0,0)=(q_7,2,R)$. بخش ۱ دومینوی $\dom{\#}{\#q_00100\#}$، بخش ۲ دومینوی $\dom{q_00}{2q_7}$ و بخش ۴ دومینوهای $\dom{0}{0},\dom{1}{1},\dom{2}{2},\dom{\sqcup}{\sqcup}$ را می‌سازند. تطابق جزئی:
\begin{ltrfig}
\renewcommand{\arraystretch}{1.3}
\begin{tabular}{l|c|c|c|c|c|c|}
top & $\#$ & $q_00$ & $1$ & $0$ & $0$ & $\#$\\ \hline
bottom & $\#q_00100\#$ & $2q_7$ & $1$ & $0$ & $0$ & $\#$
\end{tabular}
\end{ltrfig}
حالا رشته بالا $\#q_00100\#$ و پایین $\#q_00100\#2q_7100\#$ است؛ یعنی پایین یک پیکربندی ($C_2=2q_7100$) جلوتر است. اگر $\delta(q_5,0)=(q_9,2,L)$ بود، بخش ۳ دومینوهای $\dom{0q_50}{q_902}$، $\dom{1q_50}{q_912}$، $\dom{2q_50}{q_922}$ و $\dom{\sqcup q_50}{q_9\sqcup2}$ را می‌ساخت.
\end{example}

\textbf{بخش ۶ (پس از پذیرش):} وقتی $M$ به $\qacc$ می‌رسد، می‌خواهیم رشته بالا به پایین «برسد» تا تطابق کامل شود. برای هر $a\in\Gamma$ دومینوهای
\[\dom{a\,q_{\mathrm{acc}}}{q_{\mathrm{acc}}}\qquad\text{\rl{و}}\qquad\dom{q_{\mathrm{acc}}\,a}{q_{\mathrm{acc}}}\]
را بگذار. این دومینوها اجازه می‌دهند $\qacc$ نمادهای کنارش را یکی‌یکی «بخورد» تا فقط $\qacc$ بماند: مثلاً از $\#21\qacc02\#$ به $\#2\qacc02\#$، سپس $\#2\qacc2\#$ و \dots\ تا $\#\qacc\#$.

\textbf{بخش ۷ (تکمیل تطابق):} دومینوی $\dom{q_{\mathrm{acc}}\#\#}{\#}$ را بگذار. وقتی پایین به $\#\qacc\#$ ختم شده، این دومینو بالا و پایین را برابر می‌کند.

با این ساخت، $P'$ تطابقی که با دومینوی اول شروع شود دارد اگر و تنها اگر $M$ رشته $w$ را بپذیرد: هر تطابق باید پیکربندی‌های متوالی یک محاسبه درست $M$ را در پایین بسازد (دومینوها فقط تغییرهای مجاز اطراف هد را اجازه می‌دهند) و تنها راه برابر شدن بالا و پایین، رسیدن به $\qacc$ و استفاده از بخش‌های ۶ و ۷ است. برعکس، تاریخچه پذیرنده $M$ روی $w$ دقیقاً یک تطابق می‌سازد.

\textbf{تبدیل \lr{MPCP} به \lr{PCP}:} حالا باید شرط «شروع با دومینوی اول» را حذف کنیم. برای رشته $u=u_1u_2\cdots u_n$ تعریف کنید:
\[\star u=*u_1*u_2*\cdots*u_n,\qquad u\star=u_1*u_2*\cdots*u_n*,\qquad \star u\star=*u_1*u_2*\cdots*u_n*,\]
که $*$ نماد جدیدی است. اگر $P'=\left\{\left[\frac{t_1}{b_1}\right],\dots,\left[\frac{t_k}{b_k}\right]\right\}$، قرار دهید
\[P=\left\{\left[\frac{\star t_1}{\star b_1\star}\right],\left[\frac{\star t_1}{b_1\star}\right],\left[\frac{\star t_2}{b_2\star}\right],\dots,\left[\frac{\star t_k}{b_k\star}\right],\left[\frac{*\diamond}{\diamond}\right]\right\}.\]
در $P$ تنها دومینویی که بالا و پایینش با یک نماد یکسان ($*$) شروع می‌شود، اولی است؛ پس هر تطابق $P$ مجبور است با آن شروع شود. در بقیه دومینوها نمادهای بالا و پایین یکی در میان با $*$ جابه‌جا شده‌اند و دومینوی آخر $\left[\frac{*\diamond}{\diamond}\right]$ اجازه می‌دهد $*$ اضافه در انتهای رشته بالا هم جبران شود. با حذف $*$ها هر تطابق $P$ یک تطابق $P'$ (شروع‌شده با دومینوی اول) می‌دهد و برعکس.

پس $S$ روی $\enc{M,w}$ نمونه $P$ را می‌سازد و $R$ را روی $\enc{P}$ اجرا می‌کند و همان جواب را برمی‌گرداند. $S$ تصمیم‌گیرنده $\ATM$ است؛ تناقض.
\end{proof}

\section{کاهش نگاشتی}

حالا مفهوم کاهش را دقیق تعریف می‌کنیم. این تعریف به ما اجازه می‌دهد کاهش را دقیق‌تر به کار ببریم، مثلاً برای اثبات \textbf{تشخیص‌ناپذیری}، و در نظریه پیچیدگی هم کاربرد دارد. ایده: مسئله $A$ را با یک \textbf{تابع محاسبه‌پذیر} که نمونه‌های $A$ را به نمونه‌های $B$ تبدیل می‌کند، به $B$ کاهش می‌دهیم. به این کاهش \textbf{کاهش نگاشتی} یا \textbf{کاهش چند-به-یک} می‌گوییم.

\begin{definition}[تابع محاسبه‌پذیر]
تابع $f:\Sigma^*\to\Sigma^*$ \textbf{محاسبه‌پذیر} است اگر ماشین تورینگی مثل $M$ وجود داشته باشد که روی هر ورودی $w$ متوقف شود و در پایان فقط $f(w)$ روی نوارش باشد.
\end{definition}

مثلاً عملیات حسابی روی اعداد صحیح ($\enc{m,n}\mapsto\enc{m+n}$) یا تبدیل توصیف ماشین‌ها ($\enc{M}\mapsto\enc{M'}$ که $M'$ هرگز هد را از انتهای چپ بیرون نمی‌برد) توابع محاسبه‌پذیرند.

\begin{definition}[کاهش نگاشتی]
زبان $A$ به زبان $B$ \term{کاهش‌پذیر نگاشتی}{Mapping Reducible} است و می‌نویسیم $A\le_m B$، اگر تابع محاسبه‌پذیری مثل $f:\Sigma^*\to\Sigma^*$ وجود داشته باشد که برای هر $w$:
\[w\in A\iff f(w)\in B.\]
به $f$ \textbf{کاهش} از $A$ به $B$ می‌گوییم.
\end{definition}

\begin{ltrfig}
\begin{tikzpicture}[font=\small]
  \draw[thick,draw=jzMain] (0,0) ellipse (1.4 and 1.8); \draw[thick,draw=jzMain,fill=jzMain!8] (0,0.5) ellipse (0.9 and 0.8); \node at (0,1.0) {$A$};
  \draw[thick,draw=jzMain] (5,0) ellipse (1.4 and 1.8); \draw[thick,draw=jzMain,fill=jzMain!8] (5,0.5) ellipse (0.9 and 0.8); \node at (5,1.0) {$B$};
  \fill (0,0.3) circle (1.5pt); \fill (5,0.3) circle (1.5pt); \draw[->,thick,jzAccent] (0,0.3) to[bend left=15] node[above] {$f$} (5,0.3);
  \fill (0,-1.1) circle (1.5pt); \fill (5,-1.1) circle (1.5pt); \draw[->,thick,jzAccent] (0,-1.1) to[bend right=15] node[below] {$f$} (5,-1.1);
  \node at (0,-2.1) {$\Sigma^*$}; \node at (5,-2.1) {$\Sigma^*$};
\end{tikzpicture}
\end{ltrfig}
یعنی $f$ عضوهای $A$ را به عضوهای $B$ و غیرعضوهای $A$ را به غیرعضوهای $B$ می‌برد؛ پس آزمون عضویت در $A$ به آزمون عضویت در $B$ تبدیل می‌شود.

\begin{theorem}\label{thm:map-dec}
اگر $A\le_m B$ و $B$ تصمیم‌پذیر باشد، $A$ هم تصمیم‌پذیر است.
\end{theorem}
\begin{proof}
فرض کنید $M$ تصمیم‌گیرنده $B$ و $f$ کاهش از $A$ به $B$ باشد. تصمیم‌گیرنده $N$ برای $A$:
\begin{algo}{$N$}{روی ورودی $w$}
\begin{enumerate}
\item $f(w)$ را محاسبه کن.
\item $M$ را روی $f(w)$ اجرا کن و همان خروجی را بده.
\end{enumerate}
\end{algo}
اگر $w\in A$ آنگاه $f(w)\in B$ و $M$ می‌پذیرد؛ اگر $w\notin A$ آنگاه $f(w)\notin B$ و $M$ رد می‌کند. هر دو مرحله متوقف می‌شوند، پس $N$ تصمیم‌گیرنده $A$ است.
\end{proof}

\begin{corollary}
اگر $A\le_m B$ و $A$ تصمیم‌ناپذیر باشد، $B$ هم تصمیم‌ناپذیر است.
\end{corollary}

\begin{example}[کاهش نگاشتی از $\ATM$ به $\lng{HALT}{TM}$]
تابع محاسبه‌پذیر $f(\enc{M,w})=\enc{M',w'}$ را می‌خواهیم که
\[\enc{M,w}\in\ATM\iff\enc{M',w'}\in\lng{HALT}{TM}.\]
\begin{algo}{$F$}{روی ورودی $\enc{M,w}$}
\begin{enumerate}
\item ماشین $M'$ زیر را بساز:\\
$M'$ = «روی ورودی $x$: (۱) $M$ را روی $x$ اجرا کن. (۲) اگر $M$ پذیرفت، \textbf{بپذیر}. (۳) اگر $M$ رد کرد، وارد حلقه شو.»
\item $\enc{M',w}$ را خروجی بده.
\end{enumerate}
\end{algo}
$M'$ روی $w$ متوقف می‌شود اگر و تنها اگر $M$ رشته $w$ را بپذیرد (اگر $M$ رد کند، $M'$ حلقه می‌زند و اگر $M$ حلقه بزند، $M'$ هم حلقه می‌زند). اگر ورودی $F$ کدگذاری درستی نداشته باشد، $F$ رشته‌ای خارج از $\lng{HALT}{TM}$ را خروجی می‌دهد.
\end{example}

\begin{note}
اثبات تصمیم‌ناپذیری $\prob{PCP}$ هم در واقع دو کاهش نگاشتی است: $\ATM\le_m\prob{MPCP}\le_m\prob{PCP}$. اما همه کاهش‌های فصل قبل نگاشتی نیستند؛ مثلاً در اثبات $\lng{E}{TM}$، خروجی $R$ را برعکس کردیم. آن کاهش در واقع یک کاهش نگاشتی از $\ATM$ به $\overline{\lng{E}{TM}}$ است.
\end{note}

\begin{theorem}\label{thm:map-rec}
اگر $A\le_m B$ و $B$ تورینگ-تشخیص‌پذیر باشد، $A$ هم تورینگ-تشخیص‌پذیر است.
\end{theorem}
\begin{proof}
اثبات مثل قضیه \ref{thm:map-dec} است، با این تفاوت که $M$ و $N$ تشخیص‌دهنده‌اند: $N$ روی $w$ مقدار $f(w)$ را حساب می‌کند و $M$ را روی آن اجرا می‌کند. $N$ رشته $w$ را می‌پذیرد اگر و تنها اگر $M$ رشته $f(w)$ را بپذیرد، یعنی اگر و تنها اگر $w\in A$.
\end{proof}

\begin{corollary}
اگر $A\le_m B$ و $A$ تورینگ-تشخیص‌پذیر نباشد، $B$ هم تورینگ-تشخیص‌پذیر نیست.
\end{corollary}

دقت کنید که اگر $A\le_m B$ با تابع $f$، همان $f$ نشان می‌دهد $\overline{A}\le_m\overline{B}$ (چون $w\notin A\iff f(w)\notin B$). این نکته را در اثبات زیر به کار می‌بریم.

\begin{theorem}
$\lng{EQ}{TM}$ نه تورینگ-تشخیص‌پذیر است و نه هم-تورینگ-تشخیص‌پذیر.
\end{theorem}
\begin{proof}
\textbf{$\lng{EQ}{TM}$ تشخیص‌پذیر نیست.} نشان می‌دهیم $\ATM\le_m\overline{\lng{EQ}{TM}}$. پس $\overline{\ATM}\le_m\lng{EQ}{TM}$ و چون $\overline{\ATM}$ تشخیص‌پذیر نیست، $\lng{EQ}{TM}$ هم نیست. کاهش $f$:
\begin{algo}{$F$}{روی ورودی $\enc{M,w}$}
\begin{enumerate}
\item دو ماشین $M_1$ و $M_2$ را بساز:\\
$M_1$ = «روی هر ورودی: \textbf{رد کن}.»\\
$M_2$ = «روی هر ورودی: $M$ را روی $w$ اجرا کن. اگر پذیرفت، \textbf{بپذیر}.»
\item $\enc{M_1,M_2}$ را خروجی بده.
\end{enumerate}
\end{algo}
اگر $M$ رشته $w$ را بپذیرد، $M_2$ همه چیز را می‌پذیرد و $M_1$ و $M_2$ هم‌ارز نیستند. اگر $M$ رشته $w$ را نپذیرد، $M_2$ هیچ چیز نمی‌پذیرد و دو ماشین هم‌ارزند. پس $\enc{M,w}\in\ATM\iff\enc{M_1,M_2}\in\overline{\lng{EQ}{TM}}$.

\textbf{$\lng{EQ}{TM}$ هم-تشخیص‌پذیر نیست.} نشان می‌دهیم $\ATM\le_m\lng{EQ}{TM}$؛ پس $\overline{\ATM}\le_m\overline{\lng{EQ}{TM}}$ و $\overline{\lng{EQ}{TM}}$ تشخیص‌پذیر نیست. کاهش $g$ همان $F$ است، با این تفاوت که $M_1$ این بار \textbf{همه ورودی‌ها را می‌پذیرد}. اگر $M$ رشته $w$ را بپذیرد، $M_2$ هم همه چیز را می‌پذیرد و دو ماشین هم‌ارزند؛ وگرنه هم‌ارز نیستند. پس $\enc{M,w}\in\ATM\iff g(\enc{M,w})\in\lng{EQ}{TM}$.
\end{proof}

\begin{summary}
\begin{itemize}
\item \lr{PCP}: آیا دنباله‌ای از دومینوها هست که رشته بالا و پایین یکی شود؟ تصمیم‌ناپذیر است (تطابق = تاریخچه محاسبه پذیرنده؛ با هفت بخش دومینو و تبدیل \lr{MPCP} به \lr{PCP}).
\item $A\le_m B$: تابع محاسبه‌پذیر $f$ با $w\in A\iff f(w)\in B$.
\item اگر $A\le_m B$: تصمیم‌پذیری و تشخیص‌پذیری از $B$ به $A$ منتقل می‌شود، و تصمیم‌ناپذیری و تشخیص‌ناپذیری از $A$ به $B$.
\item $\lng{EQ}{TM}$ نه تشخیص‌پذیر است و نه هم-تشخیص‌پذیر.
\end{itemize}
\end{summary}
